% !TeX root = ../测度论笔记.tex
% 由主文档通过 \input 载入；请编译主文档。
\setcounter{chapter}{0} % 保留原书章号。
\chapter{代表性习题解答与学习检查}
\label{app:mt-exercises}

\section{集合运算与零测集}
第\ref{chap:mt-sets}章中，可取
$\mathcal A_1=\{\varnothing,\{1\},\{2,3\},S\}$、
$\mathcal A_2=\{\varnothing,\{2\},\{1,3\},S\}$。
二者都是 $\sigma$ 代数，但它们的并含 $\{1\}$ 和 $\{2\}$ 却不含 $\{1,2\}$，因此不对有限并封闭。

对第\ref{chap:mt-measures}章的可求和集合列，令 $E_m=\bigcup_{n\geq m}A_n$，则
\[
 0\leq\mu(\limsup A_n)\leq\mu(E_m)
 \leq\sum_{n\geq m}\mu(A_n)\longrightarrow0.
\]
这个证明不要求独立性，也不要求整个底空间测度有限。概率论将其称为第一 Borel--Cantelli 引理；反向命题并不由这一论证得到。

\section{幂函数、尖峰与积分尾部}
对 $\alpha>0$、$1\leq p<\infty$，
\[
 \int_0^1|x^{-\alpha}|^p\,\mathrm dx
 =\int_0^1x^{-\alpha p}\,\mathrm dx<\infty
 \quad\Longleftrightarrow\quad \alpha p<1.
\]
条件成立时积分为 $1/(1-\alpha p)$；等号边界出现对数发散。因此 $x^{-1/2}\in L^1(0,1)$，但不属于 $L^2(0,1)$。相对地，数列 $(1/n)$ 属于 $\ell^2$ 而不属于 $\ell^1$。

对 $f_n=n^\alpha\mathbf1_{(0,1/n)}$，有
$\|f_n\|_p^p=n^{\alpha p-1}$，故 $L^p$ 收敛到零当且仅当 $\alpha<1/p$。不论 $\alpha$ 为何，对每个固定 $x>0$，$f_n(x)$ 最终都为零；这再次区分了逐点行为与总体积分误差。

若 $f\in L^1$，函数 $|f|\mathbf1_{\{|f|>M\}}$ 几乎处处趋于零，并被 $|f|$ 控制，因此积分尾部趋零。这一结论还给出积分的绝对连续性：给定 $\varepsilon>0$，先选 $M>0$ 使尾部积分小于 $\varepsilon/2$，再令 $\mu(A)<\varepsilon/(2M)$，便有
\[
 \int_A|f|\,\mathrm d\mu
 \leq M\mu(A)+\int_{\{|f|>M\}}|f|\,\mathrm d\mu<\varepsilon.
\]
这里先选截断高度，再选集合测度，顺序不能倒置。

\section{极限唯一性与分层区间}
若 $f_n$ 同时依测度收敛到 $f,g$，则对任意 $\varepsilon>0$，
\[
 \{|f-g|>\varepsilon\}
 \subseteq\{|f-f_n|>\varepsilon/2\}\cup\{|f_n-g|>\varepsilon/2\}.
\]
右侧测度趋零，左侧不随 $n$ 变化，所以左侧测度为零。再对 $\varepsilon=1/m$ 取可数并，得到 $f=g$ 几乎处处。

第\ref{chap:mt-modes}章的分层区间示性函数满足 $\|f_{2^k+j}\|_p^p=2^{-k}$，所以对每个有限 $p\geq1$ 都在 $L^p$ 中趋于零，而每点仍在零、一之间无穷次变化。这不与控制收敛冲突：控制收敛给出从几乎处处收敛到积分收敛的充分条件，不断言反向蕴含。

\section{有限分割与不独立但不相关的例子}
若 $\mathcal G$ 由有限分割生成，每个 $\mathcal G$ 可测实函数必须在各分组上为常数。设这些常数为 $c_j$，则条件期望的积分约束为
$c_j\mathbb P(A_j)=\mathbb E[X\mathbf1_{A_j}]$。对正概率分组可以相除；零概率分组的常数可以任取有限值，例如取零，而不影响条件期望等价类。

当 $X$ 在 $\{-1,0,1\}$ 上均匀取值且 $Y=X^2$，有 $\mathbb E[X]=\mathbb E[X^3]=0$，故乘积期望分解成立。然而
$\mathbb P(X=0,Y=0)=1/3$，而 $\mathbb P(X=0)\mathbb P(Y=0)=1/9$，所以不独立。一个乘积矩只提供一个数值约束，独立性则要求所有 Borel 事件组合均满足乘法关系。

\section{进一步学习前应能说明的问题}
完成基础部分后，应能解释为什么测度只要求可数可加，为什么从上连续性需要有限性条件，为什么可测函数可以由简单函数逼近，以及积分为什么忽略可测零测集上的修改。完成收敛部分后，应能为交换极限与积分写出具体条件，并用反例区分几乎处处、依测度与 $L^p$ 收敛。

进入概率论前，应能从分布测度推导期望公式，说明密度相对于哪个测度定义，并用子 $\sigma$ 代数和积分恒等式定义条件期望。进入泛函分析前，应能解释 $L^p$ 为什么需要取等价类、范数不等式怎样使用，以及完备性证明中“选择差可求和的子列”发挥了什么作用。在本书介绍的正则性、Lusin 定理、微分定理及弱收敛基础上，后续可以进一步学习一般拓扑空间上的正则测度、覆盖与极大函数估计、$L^p$ 对偶理论以及 Sobolev 空间；这些进阶内容不必在第一遍学习时同时掌握。

\section{综合例题：一列尖峰同时检验多个概念}
\begin{example}
在 $(0,1)$ 上令 $f_n(x)=n^\alpha\mathbf1_{(0,n^{-\beta})}(x)$，其中 $\beta>0$、$\alpha\in\mathbb R$。比较它的逐点、依测度与 $L^p$ 收敛，并判断函数族何时一致可积。
\end{example}
\begin{solve}
固定 $x>0$，由于 $n^{-\beta}\to0$，最终 $x$ 不属于支撑，因此不论 $\alpha$ 为何，$f_n(x)\to0$。对任意固定阈值，坏集合的测度不超过 $n^{-\beta}$，也趋于零，所以始终依测度收敛。对有限 $p\geq1$，
\[
 \|f_n\|_p^p=n^{\alpha p-\beta}.
\]
因此 $L^p$ 收敛到零当且仅当 $\alpha p<\beta$；等号时范数保持为一，大于时范数发散。若 $\alpha>0$，本质上确界为 $n^\alpha$ 并发散；若 $\alpha=0$，它恒为一；若 $\alpha<0$，则在 $L^\infty$ 中也收敛到零。

对一致可积性，已经知道依测度收敛，因此可应用 Vitali 定理：它等价于 $L^1$ 收敛，即 $\alpha<\beta$。也可以直接检查尾部。若 $\alpha\leq0$，函数统一被常数一控制；若 $0<\alpha<\beta$，超过高度 $M$ 的项要求 $n>M^{1/\alpha}$，其尾部质量为 $n^{\alpha-\beta}$，上确界随 $M$ 增大趋零。若 $\alpha=\beta$，任意高截断仍会漏掉质量一；若 $\alpha>\beta$，连 $L^1$ 有界性都失效。
\end{solve}
这一例题中的三个参数分别描述高度、支撑长度和误差惩罚指数。逐点收敛只看到支撑最终离开每个固定点，依测度收敛只看到支撑变小，而积分收敛比较的是高度增长与支撑收缩之间的竞争。解题时先分辨每种收敛关注的信息，再计算相应量，比仅凭图像“越来越尖”作判断更可靠。

\section{综合例题：一致可积但没有共同可积上界}
\begin{example}
在 $(0,1)$ 内选择两两不交的区间 $I_n$，长度为 $2^{-n}/n$，并令 $f_n=2^n\mathbf1_{I_n}$。证明 $f_n\to0$ 于 $L^1$，函数族一致可积，但不存在可积函数同时支配全部 $f_n$。
\end{example}
\begin{solve}
这些长度之和严格小于 $\sum_n2^{-n}=1$，所以可以按长度依次排列成不交区间。每项积分为 $1/n$，因此 $L^1$ 收敛到零，Vitali 定理的必要性给出一致可积。也可直接检验：只有满足 $2^n>M$ 的项对高值尾部有贡献，而这些项的积分上确界随最低可能序号增大而趋零。

若 $g\geq f_n$ 几乎处处对每个 $n$ 成立，合并全部例外零测集后，在各不交区间上都必须有相应下界，因此
$\int_0^1g\geq\sum_n\int_{I_n}f_n=\sum_n1/n=\infty$，矛盾。这里各函数分别只占一个区间，但共同上界必须同时覆盖所有区间，其累计代价无限。
\end{solve}
这证明控制收敛的共同可积上界并非 $L^1$ 收敛的必要条件，也说明一致可积性确实比单一上界更灵活。若解题时找不到控制函数，应先检查是否可以直接估计积分或利用一致可积性，而不是把找不到某种充分条件当作不收敛的证据。

\section{综合例题：密度集中与分布极限}
\begin{example}
在 $[0,1]$ 上令 $p_n(x)=(n+1)x^n$。分析它的逐点极限、积分极限和所定义概率测度的弱极限。
\end{example}
\begin{solve}
每项非负，且积分为一，因此是概率密度。对每个 $0\leq x<1$，指数衰减快于线性增长，故 $p_n(x)\to0$；端点一的异常不影响几乎处处结论。极限函数的积分却为零，因此不存在 $L^1$ 收敛到零，Scheffé 引理的“极限仍具有正确总质量”条件失败。

对任意 $0<\delta<1$，密度在远离右端点的区域上的质量为
$\int_0^{1-\delta}p_n(x)\,\mathrm dx=(1-\delta)^{n+1}\to0$。若 $\varphi$ 在 $[0,1]$ 上连续，则一致连续性保证，在足够小的端点邻域中 $|\varphi(x)-\varphi(1)|$ 很小；在其余区域，误差不超过 $2\|\varphi\|_\infty$，而该区域质量趋零。因此
$\int_0^1\varphi(x)p_n(x)\,\mathrm dx\to\varphi(1)$，
即概率测度弱收敛于 $\delta_1$。
\end{solve}
函数几乎处处趋零与分布弱收敛到一个概率测度并不矛盾。前者观察密度相对于长度的点态数值，后者观察质量对连续测试函数的作用。极限质量集中到长度为零的单点后，无法继续由普通长度密度表示，但仍然是一个完全合法的概率测度。

\section{综合例题：平方误差中的条件平均}
\begin{example}
设 $\Omega$ 被有限分割 $A_1,\ldots,A_m$ 划分，$\mathbb P(A_j)>0$，$X\in L^2$。在所有形如 $Y=\sum_jc_j\mathbf1_{A_j}$ 的预测量中，求使 $\mathbb E[(X-Y)^2]$ 最小的系数。
\end{example}
\begin{solve}
因为分组两两不交，目标可以写成
$\sum_j\int_{A_j}(X-c_j)^2\,\mathrm d\mathbb P$。
各项只依赖自己的 $c_j$，所以可逐组最小化。展开第 $j$ 项，得到二次函数
$\int_{A_j}X^2-2c_j\int_{A_j}X+c_j^2\mathbb P(A_j)$。
二次项系数为正，唯一极小点为
$c_j=\mathbb E[X\mathbf1_{A_j}]/\mathbb P(A_j)$。
这恰是有限分割上的条件期望。

也可不用求导，直接以该均值 $m_j$ 配方：
\[
 \int_{A_j}(X-c_j)^2
 =\int_{A_j}(X-m_j)^2+\mathbb P(A_j)(c_j-m_j)^2.
\]
交叉项为零，因为 $\int_{A_j}(X-m_j)=0$。这给出最优值、唯一性及偏离最优系数时的误差增量。
\end{solve}
若某个分组概率为零，目标完全不依赖该组系数，所以不能声称所有点值系数都唯一；唯一性只在几乎必然意义下成立。这个有限维例子把抽象投影定理中的可测性、正交性和等价类三项要求都转化成了可以直接计算的条件。

\section{综合例题：积分换序先检查绝对值}
\begin{example}
对 $a>0$，计算
$I(a)=\int_0^\infty\int_0^\infty e^{-x-y}\cos(axy)\,\mathrm dx\,\mathrm dy$，
并说明交换积分和参数求导各需要哪些条件。此处不要求把最终一维积分写成初等函数。
\end{example}
\begin{solve}
首先，绝对值不超过 $e^{-x-y}$，其二维积分为一，所以 Fubini 定理允许交换次序。固定 $y$，初等积分给出
$\int_0^\infty e^{-x}\cos(ayx)\,\mathrm dx=1/(1+a^2y^2)$，
故
$I(a)=\int_0^\infty e^{-y}/(1+a^2y^2)\,\mathrm dy$。
虽然原函数未必具有方便的初等表示，这个表达已经将二维计算降为一维，并立即给出 $0<I(a)\leq1$。

若对原二重积分求参数导数，被积函数导数的绝对值不超过 $xy e^{-x-y}$，其积分等于
$(\int_0^\infty xe^{-x}\,\mathrm dx)^2=1$。
因此参数微分可移入积分号，得到
$I'(a)=-\iint xy e^{-x-y}\sin(axy)\,\mathrm dx\,\mathrm dy$。
若改对一维表达式求导，则得到
$I'(a)=-2a\int_0^\infty y^2e^{-y}/(1+a^2y^2)^2\,\mathrm dy<0$。
两种表达相同，但后一种更容易看出单调性。
\end{solve}
例题说明，不同积分表达的价值可以不同：一个方便验证存在性，另一个方便读出符号，第三个可能方便做数值计算。合法变换需要条件，而选择哪一种合法表达则取决于当前要回答的问题。

\section{证明训练：定位每一条假设的用途}
阅读一个证明时，可以把假设与实际使用位置逐一对应。有限测度常用于从上连续性、常数函数可积和从依测度误差估计积分；$\sigma$ 有限常用于把问题分成可数个有限块；非负性常用于单调收敛和不受顺序影响的求和；绝对可积性常用于有符号积分的拆分与交换；完备性则用于任意零测修改仍然可测。若某条假设在自己的复述中从未出现，应检查是否遗漏了关键步骤，或者定理确实有更弱的版本。

例如，从几乎处处收敛推出依测度收敛的证明，对尾部坏集合使用从上连续性，有限测度正是在此进入；从 $L^p$ 收敛推出依测度收敛的证明只用 Markov 型估计，不需要总测度有限。两种结论看起来都在比较收敛强弱，但它们对空间条件的依赖不同。把依赖写清，比笼统地说“强收敛推出弱收敛”更准确。

再如，Radon--Nikodym 证明中的绝对连续性用于排除“正剩余质量落在零测度块上”的情况，$\sigma$ 有限性用于找到一个具有正剩余质量的有限测度块，Hahn 分解则用于从这个块中提取可以逐点增加候选密度的正区域。这些步骤各有职责，任何一个都不能仅由候选积分有上界替代。

完成习题后，除了核对最终公式，还应检查对象是否已定义、等式在哪种意义下成立、涉及的集合是否可测、积分是否允许无穷，以及极限交换所用条件是否适用于整个序列。数学写作中的专业性首先体现在这些逻辑边界清楚，而不是使用尽可能多的术语。

